\section{Conclusion}
\label{sec:conclusion}

We document the intraday price response of EUR/USD and the
Nasdaq-100 to a thirteen-month sample of $2{,}449$ unscheduled news
events from a public real-time newsfeed, with each event
independently labelled by a large language model and benchmarked
against the standard Loughran-McDonald dictionary and the
FinBERT-tone neural classifier. After uniform Benjamini-Hochberg
false discovery rate adjustment, the news events are robustly
associated with intraday magnitude that is $1.3\times$ to
$1.9\times$ the matched-baseline level on every asset-window cell
tested, and this finding is stable under winsorisation, multivariate
controls, and a pre/post-knowledge-cutoff split. The directional
content of the LLM sentiment is modest ($49$--$54\%$ hit rate
against realised direction) and below transaction costs in a
backtest, but the LLM nonetheless outperforms the Loughran-McDonald
baseline by $2$--$6$ percentage points and the FinBERT baseline by
up to $4.5$ percentage points on the primary windows. The LLM is
also the only one of the three classifiers capable of producing
asset-specific sentiment (different USD and NDX labels on the same
headline), which it does on $73.5\%$ of events.

Two methodologically substantive findings sit alongside the main
results. The pre-event drift in the fifteen minutes before each
Discord headline is comparable in magnitude to the post-event
drift, suggesting that the public publication time is not the
informational event time in this setting; we recommend that future
work on online-news event studies obtain primary-source timestamps
where possible and report the pre-event drift as a diagnostic in
any case. And the initial price reaction extends rather than
reverses over the subsequent hours, with the sign of the
$+15$-minute return agreeing with the sign of the $+4$-hour return
on $57.6\%$ of events and median absolute moves $3$--$4.5\times$
larger over the longer window.

Two practical implications follow. For research, an LLM-based
sentiment classifier provides meaningful value over the standard
dictionary and BERT-family baselines in this kind of analysis,
both in hit rate and in the asset-specific expressiveness that
single-output classifiers cannot match; the full validation of the
LLM labels nonetheless requires manual annotation, which we are
pursuing. For practitioners using this kind of feed, the magnitude
signal is more useful than the directional signal: news events are
most reliably understood as triggers for volatility expansion, with
direction available only as a weak secondary input that does not
overcome execution costs in the naive forms tested.

We make the full pipeline---including the validation harness, the
LLM sentiment cache, the dictionary and FinBERT baselines, and the
three-way comparison---publicly available in a single repository
with seeded reproducibility, in the hope that the methodology can
be re-applied to other public newsfeeds and to a wider set of
assets.
